%% kapa: exam=1 round=36 year=2025 period=1 form=A subject=economics subject_ko=경제학원론 seq=12 number=52
%% source: Q-Net 감정평가사 기출문제 (cst00309) artlSeq=5231525 file=raw/1008/5231525/2025년 제36회 감정평가사 1차시험 1교시 문제지 원본.pdf
%% conversion: kordoc 4.24.51 -> tools/gen1.py (auto); review=proofread-sonnet; flags=none
\begin{question}{52}
생산함수가 $Q=AL^{\alpha}K^{\beta}$ 일 때 기업의 비용에 관한 설명으로 옳은 것은? (단, $Q$는 수량, $A$ 는 상수, $L$ 은 노동, $K$는 자본, $\alpha ,\beta >0$ )
\begin{choices}
  \item $\alpha +\beta =1$ 이면 장기평균비용은 1로 일정하다.
  \item $\alpha +\beta <1$ 이면 장기평균비용은 장기한계비용보다 크다.
  \item $\alpha +\beta <1$ 이면 규모에 대한 수익체증(increasing returns to scale)이 된다.
  \item $\alpha +\beta >1$ 이면 생산량이 증가함에 따라 장기평균비용은 감소한다.
  \item $\alpha +\beta >1$ 이면 규모의 불경제가 발생한다.
\end{choices}
\end{question}
\answer{4}
